\section{Introducción: por qué leer este artículo}

Anthropic es el desarrollador de la serie de modelos de lenguaje grandes Claude. En el último año, colaboró con decenas de equipos de la industria para construir sistemas de Agentes basados en LLMs. Este artículo sintetiza su experiencia práctica directa; el hallazgo central es: \textbf{las implementaciones más exitosas no dependen de marcos complejos, sino que adoptan patrones de diseño simples y composables}.

El valor único de este artículo radica en lo siguiente: no se trata de un paper académico, sino de un «manual de patrones de diseño» procedente de primera línea de la industria. El texto descompone sistemáticamente los sistemas Agentic en dos grandes categorías: \textbf{Workflow} (flujo de trabajo) y \textbf{Agent} (agente autónomo), detallando cinco patrones de Workflow y principios de diseño para Agentes, además de proporcionar casos prácticos y mejores prácticas para el diseño de herramientas en el apéndice.

\begin{importantbox}{Puntos centrales del artículo}
El éxito al construir aplicaciones con LLMs no reside en crear sistemas más complejos, sino en diseñar el sistema \textbf{más adecuado} para la necesidad. Se debe comenzar con prompts simples y optimizar mediante evaluación sistemática; solo se deben introducir sistemas agentic de múltiples pasos cuando las soluciones simples resultan insuficientes.
\end{importantbox}

\subsection{Resumen de la sección}
Este texto proviene de un resumen de prácticas de ingeniería de Anthropic, cuyo argumento central es construir sistemas agentic «de lo simple a lo complejo». Comprender este artículo ayuda a tomar decisiones arquitectónicas correctas en el desarrollo práctico.
